%=========== ΛΥΣΗ - 34ο ΘΕΜΑ Δ ===========
\begin{thema}{Δ}
\begin{erwthma}
\item Η $f$ είναι παραγωγίσιμη στο $(0,+\infty)$ και για κάθε $x>0$ είναι
\begin{align*}
f'(x)
&=\left(\frac{\ln x}{x}-\frac{1}{x}+a\right)'\\
&=\frac{1-\ln x}{x^2}+\frac{1}{x^2}\\
&=\frac{2-\ln x}{x^2}.
\end{align*}
Για $x_0=1$ έχουμε
\[
f(1)=a-1
\quad\text{και}\quad
f'(1)=2.
\]
Άρα η εξίσωση της εφαπτομένης της $C_f$ στο σημείο με τετμημένη $x_0=1$ είναι
\begin{align*}
y-f(1)&=f'(1)(x-1)\\
y-(a-1)&=2(x-1),
\end{align*}
δηλαδή
\[
y=2x+a-3.
\]
Η εφαπτομένη διέρχεται από την αρχή των αξόνων $O(0,0)$, αν και μόνο αν
\[
0=a-3.
\]
Επομένως
\[
{a=3}.
\]

\item Για $a=3$ είναι
\[
f(x)=\frac{\ln x}{x}-\frac{1}{x}+3
\]
και
\[
f'(x)=\frac{2-\ln x}{x^2}.
\]
Επειδή $x^2>0$, το πρόσημο της $f'$ είναι το πρόσημο της παράστασης $2-\ln x$. Έχουμε:
\begin{itemize}
\item $f'(x)>0\Leftrightarrow 0<x<e^2$,
\item $f'(x)=0\Leftrightarrow x=e^2$,
\item $f'(x)<0\Leftrightarrow x>e^2$.
\end{itemize}
Στον παρακάτω πίνακα βλέπουμε το πρόσημο της $f'$ και τη μονοτονία της $f$:
\begin{center}
\begin{tikzpicture}
\tkzTabInit[color,lgt=2,espcl=1.8,colorC = maincolor!50!white,colorL = maincolor!20!white,colorV = maincolor!50!white]%
{$x$ / .8,$f'(x)$ /0.8,$f(x)$/1}%
{$0$,$e^2$,$+\infty$}%
\tkzTabLine{d, +, z, -, }
\tkzTabVar{D-/{-\infty},+/{3+\frac{1}{e^2}},-/{3}}
\end{tikzpicture}
\end{center}

Άρα η $f$ είναι \textbf{γνησίως αύξουσα} στο $(0,e^2]$ και \textbf{γνησίως φθίνουσα} στο $[e^2,+\infty)$. Επομένως παρουσιάζει ολικό μέγιστο στο $x_0=e^2$, με τιμή
\[
{f(e^2)=3+\frac{1}{e^2}}.
\]

Επιπλέον,
\[
\lim_{x\to0^+}f(x)=-\infty
\]
και
\[
f(1)=2>0.
\]
Επειδή η $f$ είναι συνεχής στο $(0,+\infty)$, από το Θεώρημα Ενδιάμεσων Τιμών υπάρχει μία τουλάχιστον ρίζα
\[
\rho\in(0,1).
\]
Η ρίζα αυτή είναι μοναδική στο $(0,e^2]$, διότι η $f$ είναι γνησίως αύξουσα στο διάστημα αυτό.

Τέλος, για $x\geq e^2$ είναι
\[
f(x)>\lim_{t\to+\infty}f(t)=3>0,
\]
οπότε δεν υπάρχουν άλλες ρίζες. Συνεπώς η εξίσωση $f(x)=0$ έχει ακριβώς μία ρίζα και
\[
{\rho\in(0,1)}.
\]

\item Για κάθε $x>0$ έχουμε
\begin{align*}
F'(x)
&=\left(\frac{1}{2}\ln^2x-\ln x+3x\right)'\\
&=\frac{\ln x}{x}-\frac{1}{x}+3\\
&=f(x).
\end{align*}
Επομένως η $F$ είναι παράγουσα της $f$ στο $(0,+\infty)$.

Από το προηγούμενο ερώτημα γνωρίζουμε ότι η μοναδική ρίζα $\rho$ της $f$ ανήκει στο $(0,1)$. Άρα
\[
f(x)>0
\]
για κάθε $x\in[1,e^2]$ και το ζητούμενο εμβαδόν είναι
\begin{align*}
E
&=\int_1^{e^2}f(x)\d x\\
&=\left[\frac{1}{2}\ln^2x-\ln x+3x\right]_1^{e^2}\\
&=\left(\frac{1}{2}\cdot 2^2-2+3e^2\right)-3\\
&={3(e^2-1)}.
\end{align*}

\item Θεωρούμε τη συνάρτηση
\[
h(x)=e^2\ln x-x-e^2,\quad x>0.
\]
Για κάθε $x>0$ είναι
\[
h'(x)=\frac{e^2}{x}-1=\frac{e^2-x}{x}.
\]
Άρα η $h$ είναι γνησίως αύξουσα στο $(0,e^2]$ και γνησίως φθίνουσα στο $[e^2,+\infty)$. Επομένως παρουσιάζει ολικό μέγιστο στο $x_0=e^2$, με τιμή
\[
h(e^2)=e^2\ln e^2-e^2-e^2=0.
\]
Συνεπώς, για κάθε $x>0$,
\[
h(x)\leq0,
\]
δηλαδή
\[
{e^2\ln x\leq x+e^2}.
\]
Η ισότητα ισχύει μόνο για $x=e^2$.

Για να ορίζεται η εξίσωση, πρέπει
\[
x>0
\quad\text{και}\quad
x^2-e^4+e^2>0.
\]
Από την ανισότητα που αποδείξαμε, για κάθε $t>0$ έχουμε
\begin{align*}
e^2\ln t&\leq t+e^2\\
\Leftrightarrow\quad
\frac{\ln t-1}{t}&\leq\frac{1}{e^2}\\
\Leftrightarrow\quad
f(t)&\leq3+\frac{1}{e^2}.
\end{align*}
Η ισότητα ισχύει μόνο για $t=e^2$.

Επομένως
\[
f(x)+f(x^2-e^4+e^2)
\leq6+\frac{2}{e^2}.
\]
Η δοσμένη εξίσωση ισχύει αν και μόνο αν και οι δύο όροι του αθροίσματος λαμβάνουν τη μέγιστη τιμή τους. Άρα πρέπει
\[
x=e^2
\]
και
\[
x^2-e^4+e^2=e^2.
\]
Για $x=e^2$ η δεύτερη ισότητα επαληθεύεται και τα ορίσματα της $f$ είναι θετικά. Συνεπώς η μοναδική λύση της εξίσωσης είναι
\[
{x=e^2}.
\]
\end{erwthma}
\end{thema}
%=========== ΤΕΛΟΣ ΛΥΣΗΣ - 34ο ΘΕΜΑ Δ ===========
